\section{The \bname Dataset}
\label{sec:benchmark}

To curate a dataset on which to probe LM opinions, we must tackle three challenges.
First, we must identify topics where these opinions are relevant and curate pertinent questions for them.
Next, the questions must be designed such that we can easily extract LM opinions on them---which is challenging if the questions are fully open-ended due to the breadth of possible responses.
Finally, we need a reference distribution of human opinions from representative groups to compare LMs to.
We now discuss how we can address all these challenges by leveraging public opinion surveys.

\subsection{The power of surveys}
We observe that the aforementioned challenges in studying LM opinions also arise when attempting to measure human opinions for research or policymaking. 
The primary approach for the latter currently is to use \emph{public opinion surveys}.
According to Pew Research: ``Much of what the country [US] knows about its media usage, labor and job markets, educational performance, crime victimization, and social conditions is based on data collected through polls.''
These surveys address the first of the three challenges with the help of experts, who identify topics of public interest and carefully design questions to capture the nuances of the topic. 
To tackle the difficulties associated with analyzing open-ended responses, survey designers craft the questions to be \emph{multiple-choice}.
Finally, surveys determine humans' opinions on these topics through extensive polling of the public at large. 
(A further discussion of the meticulous data collection process followed by survey designers is provided in Appendix~\ref{app:pew_collection}.)
These factors make public opinion surveys an ideal testbed to study LM opinions, and our work develops methods for querying LMs with these surveys, as well as evaluation metrics for quantifying their alignment w.r.t. human opinions.


\subsection{Our framework}
\label{sec:methodology}
We now put forth a general methodology to convert multiple-choice public opinion surveys into datasets for evaluating LM opinions.
Consider a survey with a set of questions $Q$, where a question $q$ has a set of possible answers $A(q)$.
Each question is also categorized into a set of topics (it can have multiple associated topics), such that the questions belonging to a topic $T$ (\eg, ``guns'' for Figure~\ref{fig:prompt_example}) are denoted by $Q_T$. 
As part of the survey, each question is presented to a carefully chosen pool of participants, where every individual ($h$) must select one answer $F(h, q)$.

To use this data for our study, we need to obtain the \emph{human opinion distribution} against which we can compare LMs.
For a question, we can build this distribution by aggregating the responses over a set of human respondents $H$, \ie, $D_H(q) = \sum_{h \in H} w_h F(h, q)$.
During aggregation, we can weight respondents uniformly $w_h = 1/|H|$, or if available, using weights assigned by the survey to correct sampling biases ($\sum_{h \in H} w_h = 1$).
In this work, we will consider two different sets of respondents -- all survey respondents (O) or a demographic group such as ``Democrats'' ($G$).
We use $D_\text{O}(q)$ and $D_G(q)$ to denote the associated marginal opinion distributions respectively.

\subsection{Instantiating \bname} 
We now apply this methodology to the annual \emph{``American Trends Panel''} (ATP) polls conducted by Pew research to build the \bname dataset (details in Appendix~\ref{app:survey_curation}).
Concretely, we use \nsurveys{} ATP polls, chosen to cover a range of topics such as privacy, political views, and health.
Each poll contains two key objects that we will use for our analysis: a set of multiple-choice questions (typically $\sim 100$) and answers from respondents (typically on the order of thousands) from across the US along with their demographic information (Appendix Table~\ref{tab:survey_stats}).
We use individual survey responses -- in conjunction with demographic information and participant weights -- to obtain the per-question overall $D_{
  \text{O}}(q)$ and group-level $D_G(q)$ human opinion distributions for each of 60 demographic groups (Appendix Table~\ref{tab:survey_md}).
Pew surveys often touch upon a broad range of (often overlapping) issues---both \texttt{ATP-W26} and \texttt{ATP-W92} have questions about guns.
Thus, we further aggregate the dataset questions into the 23 coarse and 40 fine-grained topic categories shown in Appendix Table~\ref{tab:topic}.  \\

\noindent Note: While our methodology is general, the \bname dataset itself is English and US-centric. Thus, our subsequent analysis is limited to the US populace and demographic groups within (see Section~\ref{sec:discussion} for a discussion).




